
\uline{par} $(E-n)$. \uline{Si} $f$ \uline{est le morphisme structural de} $P$, \uline{alors} $f^*$, \uline{défini car} $f$ \uline{est plat}, \uline{est surjectif et définit un isomorphisme} :
\[
K_\bullet(S)/J\cdot K_\bullet(S)\xrightarrow{\sim}K_\bullet(P).
\]

Soient $\underline{E}_0=(\underline{O}_S)^n$, $e_1,\ldots,e_n$ la base évidente de $\underline{E}_0$, $\underline{E}_0^{(i)}$ le sous-Module de $\underline{E}_0$ engendré par $e_1,\ldots,e_i$. Le sous-foncteur de $G$ qui à tout $S$-schéma $S'$ associe le sous-groupe des automorphismes de $\underline{E}_{0S'}$ qui laissent invariants les sous-modules $\underline{E}_{0S'}^{(i)}$ pour tout $i$ est représentable par un sous-schéma en groupes $B$ de $G$, qui est un sous-groupe de Borel de $G$. Le faisceau quotient $\underline{G}/\underline{B}$ (pour la topologie fpqc) est représentable par le schéma des drapeaux de $\underline{E}_0$, qui sera noté $X_0$. Soient $\underline{L}_0^{(i)}$ les quotients canoniques de l'image inverse de $\underline{E}_0$ sur $X_0$, $\underline{G}_0^{(i)}$ le noyau de $\underline{E}_{0X_0}\to \underline{L}_0^{(i)}$, et $\underline{F}_0^{(i)}=\underline{G}_0^{(i-1)}/\underline{G}_0^{(i)}\simeq \Ker(\underline{L}_0^{(i)}\to \underline{L}_0^{(i-1)})$.

Par ailleurs, $B$ est produit semi-direct de son radical unipotent $U$ et de son sous-tore maximal $T$ laissant les $e_i$ invariants. Soit $Y_0=G/U$. Alors $Y_0$ est un torseur de base $X_0$ sous le groupe $T_{X_0}$. Si $S'$ est un $S$-schéma, la donnée d'un $S'$-point de $Y_0$ équivaut à la donnée, sur un $S''$ fidèlement plat quasi-compact sur $S'$, d'un automorphisme de $\underline{E}_{0S''}$ modulo les unipotents de $U_{S''}$, c'est-à-dire d'un drapeau de $\underline{E}_{0S''}$, les sous-modules du drapeau étant libres, et pour tout $i$, d'une classe d'éléments $f_i^\alpha$ libres sur $\underline{O}_{S''}$ tels que $f_i^\alpha-f_i^\beta\in \underline{G}_{0S''}^{(i)}$, soit encore à la donnée d'un drapeau de $\underline{E}_{0S''}$ et pour tout $i$ d'une base de $\underline{F}_{0S''}^{(i)}$, cette donnée étant soumise à la condition de recollement sur $S''\times_{S'}S''$; par descente fpqc, cela revient à la même donnée sur $S'$. On voit alors immédiatement que $Y_0$ est le $X_0$-schéma affine défini par l'algèbre :
